God's children already belong to him, yet they await a likeness that has not been fully disclosed. First John places those two assertions beside one another. The children are real; their life is unfinished. The promise of seeing God therefore concerns both their future and what they desire now. All Saints celebrates the fulfillment of that promise in the saints while the assembled Church still asks for mercy and growth in charity.

The Gospel describes people whose circumstances might seem to contradict their blessedness. They mourn, hunger for justice, make peace and suffer persecution. Revelation depicts a multitude that has passed through affliction and stands before God. Its members have washed their garments in the Lamb's blood. The Mass sets this vision beside Jesus' promises, and its prayers ask for help from the saints who have reached the end of the journey. Their glory and the assembly's need belong to one communion.

Two questions deserve separate treatment. The first concerns purification: how does a person who already receives God's love become capable of seeing him? Augustine explains the growth of holy desire and the believer's willing dependence on grace. Thomas Aquinas distinguishes the virtues through which a person approaches happiness from the completed happiness of seeing God. The second concerns the people Christ gathers: how can they be one, and why may their work for peace bring opposition? Bede explains the numbered and innumerable assemblies in Revelation; John Chrysostom explains peacemaking through the reconciling work of the Son.

These questions give different accounts of the same Christian life. The first asks about desire and its fulfillment. The second asks about membership, reconciliation and perseverance. The difference becomes clear at the difficult words. John's believer ``sanctifieth himself,'' yet receives sonship as a gift. Revelation gives a definite number, then shows a crowd beyond counting. Jesus blesses peacemakers, then blesses those persecuted for justice.

The prayers also speak of a difference between those already with God and those still seeking salvation. The saints continue to intercede. The proper Preface names the heavenly Jerusalem as mother and speaks of pilgrimage, help and example. At Communion, the appointed antiphon selects the promises of vision, sonship and kingdom. The Prayer after Communion asks about charity's completion and the heavenly banquet. Thus the account of holiness must explain receiving help, belonging to others and hoping for a fulfillment still to come.


